\subsubsection{Connected sums and stabilization}
\label{subsec:connected-sums}

Throughout this subsection, fix a twisted categorical input \((\cC,\eta)\).

\begin{definition}[Normalized invariant]
\label{def:normalized-invariant}
For a marked 4-manifold \((X,\alpha)\), set
\[
  \widehat{\CY}_{\cC,\eta}\bigl([(X,\alpha)]\bigr)
  :=\operatorname{Dim}(\cC)^{-1}
    \CY_{\cC,\eta}\bigl([(X,\alpha)]\bigr).
\]
\end{definition}

\begin{proposition}[Marked connected-sum formula]
\label{prop:connected-sum-formula}
For marked closed oriented 4-manifolds $(X,\alpha)$ and $(Y,\beta)$,
use the canonical splitting
\[
  H_2(X\#Y;\mathbb Z)
  \cong H_2(X;\mathbb Z)\oplus H_2(Y;\mathbb Z)
\]
to regard $\alpha\oplus\beta$ as a marking on $X\#Y$.  Then
\[
  \CY_{\cC,\eta}\bigl([(X\#Y,\alpha\oplus\beta)]\bigr)
  =\operatorname{Dim}(\cC)^{-1}
   \CY_{\cC,\eta}\bigl([(X,\alpha)]\bigr)
   \CY_{\cC,\eta}\bigl([(Y,\beta)]\bigr).
\]
Equivalently,
\[
  \widehat{\CY}_{\cC,\eta}\bigl([(X\#Y,\alpha\oplus\beta)]\bigr)
  =\widehat{\CY}_{\cC,\eta}\bigl([(X,\alpha)]\bigr)
   \widehat{\CY}_{\cC,\eta}\bigl([(Y,\beta)]\bigr).
\]
\end{proposition}

\begin{proof}
Take one-ended Kirby presentations $\mathcal H_X$ and $\mathcal H_Y$, draw
their Kirby diagrams in disjoint $3$-balls, and form their connected sum
outside both diagrams.  The resulting presentation realizes the displayed
canonical splitting on $H_2$, so the cocycle representatives and hence the
character coupons separate into the $X$- and $Y$-parts.  Ordinary labels
and dotted-projector expansions separate in the same way, while the
Reshetikhin--Turaev evaluation is multiplicative under disjoint union.

The connected-sum presentation has
\[
  h_3(\mathcal H_X\#\mathcal H_Y)
  =h_3(\mathcal H_X)+h_3(\mathcal H_Y).
\]
Consequently its normalization is
\[
  \operatorname{Dim}(\cC)^{1-h_3(\mathcal H_X)-h_3(\mathcal H_Y)}
  =\operatorname{Dim}(\cC)^{-1}
   \operatorname{Dim}(\cC)^{1-h_3(\mathcal H_X)}
   \operatorname{Dim}(\cC)^{1-h_3(\mathcal H_Y)}.
\]
Multiplying the separated finite sums proves the first formula.  This is
the marked extension of the same handle calculation in
\cite[Proposition~5.14]{Tham2021}; the new coupons do not cross the
connected-sum neck.  Dividing by $\operatorname{Dim}(\cC)$ gives the
normalized formula.
\end{proof}

\begin{definition}[Hyperbolic stabilization factor]
\label{def:hyperbolic-stabilization-factor}
Put $H=S^2\times S^2$ and define
\[
  h_{\cC,\eta}:=
    \widehat{\CY}_{\cC,\eta}\bigl([(H,0)]\bigr).
\]
\end{definition}

\begin{proposition}[Marked stabilization formula]
\label{prop:marked-stabilization-formula}
Let \((X,\alpha)\) be a marked 4-manifold.  For every $r\geq0$,
\[
  \widehat{\CY}_{\cC,\eta}
    \bigl([(X\#rH,\alpha\oplus0)]\bigr)
  =h_{\cC,\eta}^{r}
    \widehat{\CY}_{\cC,\eta}\bigl([(X,\alpha)]\bigr).
\]
Moreover $h_{\cC,\eta}\neq0$.  Consequently, for every marked 4-manifold
\((M,\beta)\) and all \(r,s\geq0\), if
\[
  (X\#rH,\alpha\oplus0)
  \cong(M\#sH,\beta\oplus0),
\]
then
\[
  \widehat{\CY}_{\cC,\eta}\bigl([(X,\alpha)]\bigr)
  =\widehat{\CY}_{\cC,\eta}\bigl([(M,\beta)]\bigr)
   h_{\cC,\eta}^{s-r}.
\]
\end{proposition}

\begin{proof}
The first identity follows by induction from
\cref{prop:connected-sum-formula}.  At the zero marking the character
coupons disappear, so \cref{rem:identity-character-specialization} identifies
$h_{\cC,\eta}$, up to the displayed nonzero global-dimension
normalization, with the untwisted hyperbolic stabilization scalar.  Reutter
records the Crane--Yetter--Kauffman theory of a ribbon fusion category as an
indecomposable semisimple oriented four-dimensional TQFT
\cite[Example~2.3]{Reutter2023}; his nonvanishing theorem therefore applies
\cite[Theorem~3.5]{Reutter2023}.  Since
$\operatorname{Dim}(\cC)\neq0$, this proves $h_{\cC,\eta}\neq0$.

Apply marked diffeomorphism invariance to
the displayed stable diffeomorphism, substitute the first identity on both
sides, and cancel the nonzero powers of $h_{\cC,\eta}$.
\end{proof}
